\documentclass[11pt]{article}
\usepackage[T1]{fontenc}
\usepackage{lmodern}
\usepackage[margin=1in]{geometry}
\usepackage{booktabs}
\usepackage{amsmath}
\usepackage[hidelinks]{hyperref}

\title{Are random combinations of MLP neurons as interpretable as neurons?\\
\large Pilot study on Qwen3-0.6B with blinded agentic explanations}
\author{Rohan Buluswar}
\date{October 2026}

\begin{document}
\maketitle

\section{Question}
The essay \emph{``Neural Networks Are Already Clean''} argues that a network's own neurons are privileged, legible
building blocks. We test this against a null hypothesis: if random linear combinations of neurons are explained as
well as single neurons by the same procedure, neurons are not special. Explanations are produced and evaluated by
language-model agents that are blind to the type of unit they study.

\section{Setup}

\paragraph{Units.} At a decoder layer, let $p \in \mathbb{R}^{d_{\mathrm{ff}}}$ be the MLP hidden activation (the
input of \texttt{down\_proj}, $\mathrm{silu}(W_g x)\odot W_u x$; $d_{\mathrm{ff}}=3072$). A unit is a direction and
offset, $z = p^\top a + b$, evaluated at every token (position 0, an attention sink, is excluded everywhere). The
conditions are listed in Table~\ref{tab:conditions}. Combinations are centered and scaled to unit variance; neurons
keep their natural scale.

\begin{table}[h]
\centering\small
\begin{tabular}{ll}
\toprule
Condition & Construction \\
\midrule
\texttt{neuron} & one coordinate of $p$ (baseline) \\
\texttt{ksparse\_k}, $k\in\{2,5,20,100\}$ & $k$ random neurons, Gaussian weights on z-scored neurons \\
\texttt{related\_k}, $k\in\{5,20\}$ & a neuron and its $k-1$ most correlated neurons, sign-aligned (positive control) \\
\texttt{dense\_raw} & random direction in $\mathbb{R}^{3072}$ \\
\texttt{dense\_white} & random direction after ZCA whitening \\
\bottomrule
\end{tabular}
\caption{Unit conditions. All units of a layer are drawn at random (fixed seed), with no selection on appearance.}
\label{tab:conditions}
\end{table}

\paragraph{Model and data.} Qwen3-0.6B (28 layers), layers 7, 14 and 21. Corpus: 3M training and 0.3M held-out
tokens (FineWeb, FineWeb-Edu, code, Wikipedia in 7 languages; documents up to 1024 tokens; every tenth document held
out). One streaming pass computes activation statistics (1M tokens) to build the units; a second pass caches, per
unit, the top-1024 firings in each direction, the top-64 with activation windows, and a uniform token sample.

\paragraph{Agents.} Every agent is a separate headless Claude Code session (\texttt{claude -p}, Claude Sonnet 5.5)
started in a folder that contains only the files for one unit, with no instruction files (\texttt{CLAUDE.md}) above it;
the blinding key and ground truth are stored elsewhere. A unit is identified only by an opaque id (e.g.\
\texttt{u00102}) and its layer; agents are never told the condition, $k$, which neurons or weights make up the unit,
its raw activation scale, or its sparsity statistics. All conditions receive identical instructions.

\paragraph{What the explainer sees.} A one-page view (\texttt{view.md}, the same content as \texttt{data.json}) of
the unit's \emph{primary pole}, the direction with the more extreme corpus firings. Every activation is given as a
fraction of the unit's maximum on that pole (1.00 = strongest firing in the corpus):
\begin{itemize}
\item the 24 strongest firings, at most one per document, each as a window of 48 tokens before and 16 after the
firing token; the firing token is marked \texttt{[[tok|f]]} and every other token in the window that reaches 10\% of
the maximum is tagged \texttt{«tok|f»};
\item the 12 tokens it fires on most often and the 8 most frequent next tokens, over its 100 strongest firings;
\item up to 6 further examples in each of three strength bands (at least 50\%, 30\% and 10\% of the maximum), with
32 tokens before and 8 after;
\item 8 uniformly random tokens with their activation, as a baseline for ordinary text;
\item the 8 strongest firings of the opposite pole, with their token statistics.
\end{itemize}

\paragraph{What the explainer can do.} Agents may only read files and write files inside their own folder: the
university's Enterprise Claude account refuses every shell command and every external (MCP) tool call in headless
mode, regardless of allow-lists, so no tool that requires running code is available.
\begin{itemize}
\item \emph{In the pilot}, the folder also contained a helper script for browsing the data, which agents could not
run; sessions took at most 5 turns, so explanations rest on the one-page view. The explainer also wrote the test texts
(5 predicted-firing, 5 near-miss, 5 unrelated) used for the rate metrics. For units whose primary pole is negative,
the strength labels of the top examples were shown with the wrong sign (Section~\ref{sec:problems}).
\item \emph{Current design.} The helper script is replaced by \texttt{long\_contexts.md}: 384 tokens of context before
and 32 after each of the 12 strongest firings. Optionally, the explainer can query the model: it writes up to 10 texts
to \texttt{queries/001.txt}, and a relay outside its control runs them through the model and returns
\texttt{queries/001.out.md}, showing the unit's activation on every token in the same format as the view (texts are
cut at 256 tokens, which the answer says; the first token is never scored). The budget is 20 texts per unit, the same
for every condition, and the relay only ever queries the unit of the folder it reads from. The explainer no longer
writes test texts.
\end{itemize}
The explainer's output is an explanation ($\le$60 words) of what the primary pole responds to, where in the pattern it
fires ($\le$20 words), a list of independent triggers (\emph{parts}, 1--8), an optional remark on the other pole, and a
confidence (1--5). Limits: 40 turns and 25 minutes per session.

\paragraph{Explanation-only agents.} Both see the explanation and the ``fires on'' note and nothing else.
\begin{itemize}
\item \textbf{Writer} (in the pilot: \emph{blind-max writer}): 3 texts meant to make the unit fire as strongly as
possible. In the current design it also writes the 5 predicted-firing, 5 near-miss and 5 unrelated test texts, so all
test texts measure what the explanation conveys rather than what the explainer saw.
\item \textbf{Simulator}: rates 20 held-out contexts 0--10 for activation on a marked token. Ten are spread over the
unit's 60 strongest held-out firings (group \emph{top}) and ten are uniformly random tokens (group \emph{random}),
shuffled together.
\end{itemize}

\paragraph{Pilot size and cost.} 8 units per condition and layer were explained (216 units; 24 per condition). Agent
cost: \$39.26 (\$0.18 per unit over the three roles), with no failed sessions.

\paragraph{Metrics.} \emph{Simulator correlation}: Spearman correlation between simulator ratings and true
activations. \emph{Rate} metrics: share of a unit's probe texts that ``fire strongly''; \emph{margin} =
rate(predicted) $-$ rate(near-miss). \emph{Blind-max ratio}: best blind-max text's activation divided by the unit's
corpus maximum. Comparisons to \texttt{neuron} use a layer-stratified bootstrap 90\% interval of the difference in
means (TOST-style verdicts with margin 0.10 for correlations and rates).

\paragraph{What ``fires strongly'' means.} The original rule, at least 30\% of the unit's corpus maximum on some
token, is not comparable across conditions: dense directions are nearly Gaussian, so 30\% of their maximum is close to
ordinary noise (Table~\ref{tab:main}, last two columns). We therefore added a \emph{per-unit} threshold: a text fires
if its maximum exceeds the 95th percentile of maxima over 500 random held-out windows of the same length, computed
separately for each unit. Both definitions are reported.

\section{Results}

\begin{table}[h]
\centering\small
\begin{tabular}{lccccc}
\toprule
& Simulator & & & \multicolumn{2}{c}{Random-text firing rate} \\
\cmidrule(lr){5-6}
Condition & correlation & Parts & Confidence & 30\% of max & per-unit 95th pct \\
\midrule
\texttt{neuron}       & \textbf{0.66} & 2.7 & 3.3 & 0.02 & 0.05 \\
\texttt{related\_5}   & \textbf{0.71} & 2.6 & 3.3 & 0.06 & 0.07 \\
\texttt{related\_20}  & 0.50 & 2.5 & 3.1 & 0.05 & 0.06 \\
\texttt{ksparse\_2}   & 0.57 & 3.1 & 3.3 & 0.02 & 0.08 \\
\texttt{ksparse\_5}   & 0.58 & 2.9 & 3.2 & 0.01 & 0.09 \\
\texttt{ksparse\_20}  & 0.51 & 3.1 & 3.1 & 0.01 & 0.04 \\
\texttt{ksparse\_100} & 0.39 & 3.3 & 2.9 & 0.02 & 0.04 \\
\texttt{dense\_white} & 0.32 & 3.6 & 2.7 & 0.20 & 0.06 \\
\texttt{dense\_raw}   & 0.49 & 3.4 & 3.1 & 0.58 & 0.07 \\
\bottomrule
\end{tabular}
\caption{Per-condition means over 24 units (three layers pooled). The last two columns check the threshold
definitions: random text should rarely ``fire'' under either rule, equally across conditions.}
\label{tab:main}
\end{table}

\paragraph{The simulator separates neurons from random combinations.} Simulator correlation decreases with the number
of mixed neurons ($0.66$ for neurons; $0.57, 0.58, 0.51, 0.39$ for $k=2,5,20,100$; $0.32$ for whitened dense
directions), while the positive control \texttt{related\_5} matches neurons. Differences from neurons (90\% intervals):
\texttt{ksparse\_100} $-0.26\,[-0.36,-0.16]$ and \texttt{dense\_white} $-0.34\,[-0.43,-0.24]$ are clearly worse than
the 0.10 margin; \texttt{dense\_raw} $-0.16\,[-0.26,-0.07]$, \texttt{ksparse\_20} $-0.15\,[-0.25,-0.05]$ and
\texttt{related\_20} $-0.16\,[-0.27,-0.05]$ are below neurons but inconclusive against the margin. Explanations of
mixtures also list more independent triggers and come with lower confidence.

\paragraph{The gap is largest in the late layer.} Table~\ref{tab:layers} splits the simulator correlation by layer
(8 units per cell). At layer 21, six of eight conditions are significantly below neurons (all but \texttt{related\_5}
and \texttt{ksparse\_2}), and the decline with $k$ is steep. Layer 14 is the hardest for every condition: its neurons are the least sparse (median excess kurtosis 17,
versus 102 and 479 at layers 7 and 21).

\begin{table}[h]
\centering\small
\begin{tabular}{lcccc}
\toprule
Condition & Layer 7 & Layer 14 & Layer 21 & All \\
\midrule
\texttt{neuron}       & 0.59 & 0.59 & 0.78 & 0.66 \\
\texttt{related\_5}   & 0.65 & 0.69 & 0.80 & 0.71 \\
\texttt{related\_20}  & 0.49 & 0.41 & 0.60$^*$ & 0.50 \\
\texttt{ksparse\_2}   & 0.48 & 0.55 & 0.69 & 0.57 \\
\texttt{ksparse\_5}   & 0.51 & 0.55 & 0.69$^*$ & 0.58 \\
\texttt{ksparse\_20}  & 0.49 & 0.47 & 0.56$^*$ & 0.51 \\
\texttt{ksparse\_100} & 0.47 & 0.41$^*$ & 0.30$^*$ & 0.39 \\
\texttt{dense\_white} & 0.31$^*$ & 0.51 & 0.14$^*$ & 0.32 \\
\texttt{dense\_raw}   & 0.46 & 0.43$^*$ & 0.59$^*$ & 0.49 \\
\bottomrule
\end{tabular}
\caption{Simulator correlation by layer. $^*$: 90\% bootstrap interval of the difference from neurons in the same
layer excludes zero.}
\label{tab:layers}
\end{table}

\paragraph{What drives the simulator score.} Splitting the 20 items: within the 10 random items the correlation is
near zero for every condition (true activations there are essentially zero). The score is mostly the separation of
strong from random items (AUC 0.84 for neurons, 0.90 for \texttt{related\_5}, 0.67--0.69 for \texttt{dense\_white} and
\texttt{ksparse\_100}). Part of the drop is false alarms: the simulator rated 71\% of random items 0 for neurons but only
25\% for \texttt{dense\_white}. Within the 10 strong items, neurons are the most predictable (0.45 versus 0.23--0.33),
i.e.\ their explanations also capture firing intensity.

\paragraph{The text-based metrics do not discriminate.} The explainers' own predicted-firing texts succeed at similar rates in every
condition (88--97\% under the per-unit threshold, 71--93\% under the 30\% rule), and the median blind-max ratio is 0.71--0.90 in every condition (four
mixtures are formally equivalent to neurons). Explanations quote the strongest examples, so writers can reproduce
them without the explanation capturing a concept.

\section{Problems found and fixed}\label{sec:problems}
\begin{itemize}
\item \textbf{Threshold bias}: the 30\%-of-maximum rule counts 58\% of random texts as firing for \texttt{dense\_raw}
(2\% for neurons). Fixed by the per-unit calibrated threshold (4--9\% for every condition on independent random text).
\item \textbf{Sign error} for units whose primary direction is negative (about half): the simulator's answer key was
negated, hiding the effect (correlations $+0.52$ and $-0.53$ by direction averaged to zero). Fixed and rescored
without new agent runs. The same error showed negative strength labels in those units' example lists; scores of the
affected units do not differ from the others.
\item \textbf{Agent permissions}: the CLI refused agents' writes in headless mode until edits were explicitly allowed
inside their own folder. Shell commands and external tool calls stay refused on this account even when allowed by
name, so the pilot's helper script was unusable; it was replaced by pre-rendered files, and model queries go through
files answered by a relay.
\end{itemize}

\section{Limitations}
Small samples (24 units per condition, 20 simulator items per unit); a 0.6B model, whose neurons the essay expects to be
less clean than in larger models; analysis choices were made after seeing this run, so it is a calibration study, not a
pre-registered test; explainer, writer and simulator are the same model; explanations quote top examples, which the
writer can copy. Agents cannot run code, but whether reading files outside their folder is refused was not tested;
the blinding key is therefore kept on a different machine from the agents in the cluster setup.

\section{Changes since the pilot, and next step}
The pipeline now has the design described above: the explainer only explains (optionally with 20 model queries), an
explanation-only writer produces all test texts, all texts are scored in one model pass, and simulator results are
also reported per item group (correlation within top items, within random items, top-versus-random AUC, and mean
rating on random items, i.e.\ false alarms). A trial on 4 pilot units (two neurons, a whitened dense direction and a
\texttt{ksparse\_100} mixture) ran end to end with real agents: 3 of 4 explainers queried the model (5--6 texts each,
in one round), used the answers to confirm or narrow their hypotheses, and the writer's test texts separated neurons
from mixtures where the explainer's had not. Single units vary strongly between runs, so 4 units cannot show whether
querying improves explanations.

Next: Qwen3-8B on the cluster, layers 9, 18, 27 of 36 ($d_{\mathrm{ff}}=12288$), the full 50M-token corpus with
2048-token documents, 200 units per condition and layer cached, and 24 per condition and layer explained (648 units).
Before running agents, fix the primary metric (simulator correlation), the threshold (per-unit 95th percentile), the
equivalence margins and whether queries are enabled. Estimated cost: roughly 3--8 GPU-hours on one L40S or A100, and
about \$130--170 of agent usage at the measured \$0.20--0.25 per unit.

\end{document}
